%%
% 第一章：绪论
%%

\chapter{绪论}
\label{chap:introduction}

\section{研究背景与意义}

\subsection{研究背景}

面向非结构化环境的机械臂操作是机器人学与人工智能交叉研究的前沿方向。近年来，视觉--语言--动作（Vision-Language-Action, VLA）模型将多模态大模型的语义理解与机器人动作生成相统一，
为"自然语言条件下的灵巧操作"提供了新的技术路径\cite{black2024pi0,black2025pi05}。
与此同时，工业界与学术界相继开源了面向模仿学习与数据标准化的软件栈（如 Hugging Face LeRobot），
降低了数据采集与策略训练的工程门槛。然而，从公开算法到具体实验室的异构硬件（高精度力控臂、多路相机、实时控制链路），
中间仍普遍存在接口割裂、时序对齐困难、示教可重复性差等问题；尤其是面向 $\pi_{0.5}$ 等新一代分层 VLA 的部署，
往往要求稳定的底层执行、与通用数据集规范对齐的观测--动作字段、以及可诊断的在线闭环，这些能力难以仅靠"更换若干超参数"获得。

在此背景下，领域重要性体现在两方面。首先，桌面抓取与放置类任务在物流分拣、
实验室自动化与服务业场景中需求广泛，且同时考验全局场景理解与接触阶段的精细反馈；
其次，VLA 的论文级结果多在特定机器人形态与数据规模下呈现，若缺少与本平台一致的全栈贯通，
则难以判断性能瓶颈究竟来自模型、数据还是部署方式。
本实验室在引进 弗兰卡 等先进臂式平台并探索 VLA 软硬件闭环的过程中，围绕实时控制、开源框架扩展与策略实装等关键环节开展系统性建设。本文作为该建设链条中"桌面抓取"方向的阶段性工作，独立完成了从硬件接入、示教采集、模型训练到实机闭环评测的全栈技术链路。作者与课题组其他成员在平台搭建与数据采集等环节存在合理分工协作，但本文所涉及的系统设计、接口实现、训练方案与实验分析均由作者主导完成。

\subsection{研究问题与研究目标}

研究问题。本文围绕 VLA 机械臂桌面操作的全栈落地，重点探讨三个层次的问题。首先，在系统集成层面，如何将 弗兰卡 实时控制与 LeRobot 通用机器人/遥操作抽象对接，使观测--动作模式定义、录制流程与后续训练入口保持一致？其次，在数据层面，如何组织面向桌面抓取的多相机示教数据，使其符合 LeRobot v2.1 规范并支撑 $\pi_{0.5}$ 微调？最后，在算法与部署层面，在有限样本条件下，$\pi_{0.5}$ 经迁移训练后能否完成静平台单臂抓取--搬运--放置，以及在线执行策略如何影响成功率与失败恢复？

研究目标。本文旨在构建"底层执行—数据采集—模型训练—实机部署"的闭环系统；以苹果抓取并放置至指定区域为任务实例，完成可复现的接口扩展、数据集构建、$\pi_{0.5}$ 训练配置与对比实验，并形成可推广的部署经验。

研究假设。本文提出两项研究假设。假设 H1（迁移可行性）认为，在加载 $\pi_{0.5}$ 预训练权重的前提下，规模在数十条轨迹量级的 LeRobot 规范数据集足以支撑桌面窄分布任务的微调，使策略学到稳定的粗定位—对准—夹取—搬运行为模式。假设 H2（闭环敏感性）认为，对于易滑移、强接触的桌面物体，短时域动作块执行与高频重推理的部署方式相比偏异步、长开环的执行方式，将显著提升抓取成功率与过程稳定性；性能差异主要归因于闭环方式而非单一权重差异。

\section{国内外研究现状}

\subsection{视觉—语言对齐与空间推理}

在 VLA 之前就已有大量工作探索如何将语言与可操作视觉表征结合。
Cliport\cite{shridhar2021cliport} 将 CLIP 式语义与 ``where'' 空间变换网络结合，实现语言条件下的桌面抓取；
SayCan\cite{ahn2022saycan} 用语言模型产生子任务，再用价值函数筛选当前环境下可执行的一步，缓解 ``会说不会做''；
VIMA\cite{jiang2022vima} 则支持图文交错的多模态提示，为 ``指图操作'' 类交互奠定基础。
这一阶段的工作说明：仅有语言或仅有像素都不足以应对长程操作，跨模态对齐与可执行性约束必须协同设计。

在空间推理与三维世界建模方面，
SpatialVLM\cite{chen2024spatialvlm} 赋予 VLM 度量级空间判断能力；
3D-VLA\cite{zhen20243dvla} 将动作生成与三维世界建模耦合，面向长程后果预测；
Surfer\cite{ren2025surfer} 则以世界模型显式预测场景状态转移，为语言条件操作中的动作选择与物理后果建模提供了另一类思路；
OK-Robot\cite{liu2024okrobot} 系统分析了 ``堆叠开放知识模块'' 时真正决定成败的工程要素，
对无再训练或少再训练的快速部署具有参考价值。

\subsection{具身基础模型、开源VLA与动作生成}

PaLM-E\cite{driess2023palme} 将机器人感知流注入超大语言模型，实现长程规划与场景理解；
RoboCat\cite{bousmalis2023robocat} 强调自改进与跨机械臂形态迁移；
体素化动作建模方面，Perceiver-Actor\cite{shridhar2022perceiveractor} 在三维占用栅格上输出动作，有利于遮挡与多物体堆叠。
当前社区主流进一步转向以真机机器人数据与统一数据模式定义 为约束的高容量策略网络，而非仅在仿真或单一形态上验证的通才演示。
国内关于多模态大模型具身智能体的综述也将具身感知、任务规划与低层动作控制作为关键组成，说明 VLA 正在从单点模型演进为覆盖感知、规划和执行的系统方法\cite{zhao2025multimodalEmbodiedAgentsReview}。

在从大规模预训练到开源 VLA 的演进中，
RT-1\cite{brohan2022rt1} 证明在超大规模真机数据上训练的高容量模型可获得鲁棒策略；
RT-2\cite{brohan2023rt2} 将网络图文先验与离散动作 token 结合，使 ``网络知识'' 可迁移到机器人控制；
Open~X-Embodiment\cite{oneill2024openxembodiment} 联合多实验室、多本体数据训练 RT-X，标志着跨硬件统计迁移成为社区共识。
开源侧，OpenVLA\cite{kim2024openvla} 与 Octo\cite{ghosh2024octo} 提供了可本地微调、可替换相机与语言条件的通用策略；
RoboFlamingo-Plus\cite{wang2025roboflamingoplus} 则进一步融合深度与 RGB，缓解纯二维投影造成的深度歧义。

在仿真、数据编排与生成式 VLA 方向，
RoboCasa\cite{nasiriany2024robocasa}、ManiSkill3\cite{taole2024maniskill3} 等强调大规模仿真与 GPU 并行，用于弥补真实示教成本；
AutoRT\cite{ahn2024autort} 展示如何用具身大模型编排多台机台采集；
GR-2\cite{cheang2024gr2} 探索视频—语言先验向机器人动作的生成式迁移。
扩散模型在具身模仿学习中的应用也受到系统总结，其优势在于能够表达多峰动作分布，并通过迭代去噪机制生成平滑、稳定的操作轨迹\cite{li2026diffusionEmbodiedImitationLearningReview}。
在学习机制层面，$\pi^*_{0.6}$\cite{physicalintelligence2025pi06} 引入经验学习与修正闭环，
FAST\cite{pertsch2025fast} 与知识绝缘\cite{driess2025knowledge} 则从动作分词与训练效率角度推进 VLA。

在连续动作建模方面，
$\pi_0$\cite{black2024pi0}、$\pi_{0.5}$\cite{black2025pi05} 将预训练视觉—语言主干
（如 PaliGemma\allowbreak\cite{beyer2024paligemma}）与动作专家结合，
并以 flow matching\cite{lipman2022flowmatching} 建模连续动作分布，
代表当前面向真机高频控制的一条主流路线。
本文即以 $\pi_{0.5}$ 为主线，在其开放权重与训练栈上进行任务迁移。

\subsection{遥操作、桌面操作与部署共识}

Mobile~ALOHA\cite{fu2024mobilealoha} 与 ALOHA~Unleashed\cite{zhao2024alohaunleashed} 推动低成本的全身示教与复杂双手技能；
APT\cite{wen2023apt} 用任意点轨迹统一可形变物体与刚体操作，为动作表示提供新视角。
这些数据范式与 LeRobot 的统一张量规范相衔接，使社区能够快速沉淀可训练数据集。

在桌面操作与部署侧，
桌面抓取在几何约束、摩擦与滚动等非结构扰动下对接触瞬间的闭环修正极为敏感。
既有研究既关注抓取规划与学习型策略本身，也日益重视推理频率、动作块执行方式与异步/同步控制对成功率的影响；
但在公开讨论中，同一训练权重在不同部署脚本下的巨大性能落差仍相对缺乏系统记录，尤其对 VLA 动作块输出而言，部署策略往往与模型结构同等关键。
此外，抓取后的物理属性变化与载荷适应也会影响闭环控制稳定性，例如 FlyAware\cite{ye2026flyaware} 通过视觉估计与抓后适应处理空中操作中的惯性变化，提示接触后的在线校正同样是机器人操作系统的重要环节。

现有研究的不足与研究空白。
综合上述脉络，现有工作尚存在以下与本文课题密切相关的研究空白。

首先，平台落地鸿沟。公开 VLA 与 LeRobot 生态多以若干 ``标准型'' 机器人为第一公民，
弗兰卡 等平台若缺少与 Robot/Teleoperator 约定一致的实现，则难以直接复用社区训练脚本，
更无法保证示教数据与策略消费端的字段一一对应；而跨本体成果\cite{oneill2024openxembodiment,kim2024openvla,ghosh2024octo}虽强，
仍不能自动消解单平台线缆级工程问题。

其次，示教质量与可复现性。低成本遥操作\cite{fu2024mobilealoha,zhao2024alohaunleashed}与灵活轨迹表示\cite{wen2023apt}为数据提供了方法学支撑，
但安装误差、轴向反向与夹爪量程不一致会导致关节空间系统偏差；若缺少可文件化、可自动解算的标定流程，
则小规模高质量数据难以稳定获得，仿真合成数据\cite{nasiriany2024robocasa,taole2024maniskill3}亦无法替代本机一致性。

再次，小规模专用数据下的迁移验证。$\pi_{0.5}$\cite{black2025pi05} 原论文面向移动双臂与高覆盖多相机配置，
而静态单臂桌面任务的观测—动作维度更紧凑；如何在十余小时以下、数十条轨迹量级完成可用的任务特化微调，
并与实机接口严格对齐，仍需面向具体平台给出可检验的实现路径——这与开放世界报告\cite{black2025pi05,physicalintelligence2025pi06}形成互补问题。

最后，部署闭环与性能解释。动作块预测\cite{pertsch2025fast}在理论上利于平滑与快速推理，但若在线执行仍偏长时开环，
接触敏感任务可能出现成功率坍塌；现有文献对 ``同一模型、不同闭环脚本'' 的对比往往不足，
且空间/三维推理增强\cite{chen2024spatialvlm,zhen20243dvla,ren2025surfer}与开放系统集成经验\cite{liu2024okrobot}很少落实到弗兰卡+双相机+VLA 权重这一典型教学实验配置上。

由此，本文将问题凝练为：在实验室自建 弗兰卡--LeRobot--OpenPI（$\pi_{0.5}$）全流程的前提下，能否以标准化示教数据与针对性接口适配，在桌面抓取任务上实现可用的 VLA 策略，并厘清部署闭环对实机性能的决定性作用？

\section{本文组织架构}

\subsection{研究方法与实验环境}

本节概述全文采用的工程与研究思路，并给出实验场所软硬件条件及拟开展的实验内容；其中软硬件配置汇总见表~\ref{tab:intro-lab-env}。

（1）硬件环境。
机械臂：Franka FR3 协作机械臂（7-DoF，重复定位精度约 $\pm 0.1\,\mathrm{mm}$），末端配备知行 CTAG2F90C 双指夹爪。
感知设备：腕部末端安装 Intel RealSense D445 深度相机（RGB 分辨率 $640\times480@30\,\mathrm{Hz}$，深度量程约 $0.1$--$10\,\mathrm{m}$），获取 手眼视角；若配置固定场景机位，则与腕部相机形成互补双视角，外参与同步策略在第二章统一约定。
计算设备：实验室配备 NVIDIA RTX A6000 显卡（4 张）、Intel Xeon Platinum 8336C 处理器（2 路，32 核 64 线程/路）、DDR4 内存（8 条，每条 $32\,\mathrm{GB}$），用于模型训练与高频推理实验。
桌面场景：标准办公桌（$120\,\mathrm{cm}\times 60\,\mathrm{cm}$），摆放笔具、积木、文件、杯子等常见物件；辅以可调照度 LED 光源（约 $500$--$2000\,\mathrm{lux}$），便于考察光照变化对感知的影响。

（2）软件环境。
操作系统采用 Ubuntu 22.04 LTS。
深度学习侧以 PyTorch 2.1 为主进行模型训练；机械臂实时控制沿用第二章所述基于 libfranka 与轨迹平滑的实时控制服务，上位机经套接字协议与 LeRobot 客户端协同。
机器人数据与遥控录制依托 LeRobot SDK（v2.1 数据布局）；若需与导航、多传感器节点联动，可在同一操作系统上按需启用 ROS~2 相关接口封装作为辅助中间件。
数据处理与可视化采用 NumPy、Matplotlib，抽检与调试辅以 Rerun 等工具。

\begin{table}[htbp]
  \centering
  \small
  \caption{实验平台软硬件条件一览}
  \label{tab:intro-lab-env}
  \begin{tabular}{>{\raggedright\arraybackslash}p{1.6cm}>{\raggedright\arraybackslash}p{2.8cm}>{\raggedright\arraybackslash}p{8.4cm}}
    \toprule
    类别 & 项目 & 配置说明 \\
    \midrule
    \multirow{5}{*}{硬件} & 机械臂与夹爪 & Franka FR3（7-DoF，重复定位精度约 $\pm 0.1\,\mathrm{mm}$）；知行 CTAG2F90C 双指夹爪 \\
      & 视觉传感器 & Intel RealSense D445（腕部 手眼；RGB $640\times480@30\,\mathrm{Hz}$；深度约 $0.1$--$10\,\mathrm{m}$）；可选场景固定机位构成双视角 \\
      & 计算节点 & RTX A6000 $\times 4$；Xeon Platinum 8336C $\times 2$（32C64T/路）；DDR4 $32\,\mathrm{GB}\times 8$ \\
      & 场景与光照 & 办公桌 $120\times 60\,\mathrm{cm}$；典型桌面杂物； \\
      & 网络与控制 & 局域网连接 RT PC（运行实时控制服务）与训练/推理工作站；千兆以太网等 \\
    \midrule
    \multirow{4}{*}{软件} & 系统与运行时 & Ubuntu 22.04 LTS；Python 虚拟环境与依赖锁定 \\
      & 训练框架 & PyTorch 2.1；OpenPI/$\pi_{0.5}$ 训练配置（见第三章） \\
      & 机器人与数据 & LeRobot SDK；实时控制服务（libfranka）；可选 ROS~2 栈 \\
      & 辅助工具 & NumPy；Matplotlib；Rerun（录制质量抽检） \\
    \bottomrule
  \end{tabular}
\end{table}

\subsection{实验方案与技术路线}

实验内容初步归纳为三个方面。首先是桌面操作数据集构建：预定采集覆盖多类物体（如苹果、正方体积木等），规模不少于 30 条有效轨迹，每条轨迹包含约 $200$--$500$ 帧多模态数据，严格遵循 LeRobot v2.1 规范并在采集管线侧通过时间戳对齐。其次是模型训练与微调：在 $\pi_{0.5}$ 开源权重基础上，将第二章构建的数据集接入 OpenPI 训练接口，采用任务特化微调策略，并在训练过程中监控损失及离线指标。最后是系统测试与性能评估：在若干代表性场景下评测抓取成功率、末端定位误差及指令跟随正确率，并对比遮挡、照度变化等条件下的性能衰减，分析感知与动作生成链条中的薄弱环节。

以上三方面内容分别对应第二章（系统集成与数据采集）和第三章（模型训练与实机评测），具体方案与指标见表~\ref{tab:design-tasks}。

\subsection{设计任务与指标体系}

为将"系统集成—数据采集—模型训练—实机评测"的全栈目标落实为可检验的具体任务，本文将毕设工作分解为四项设计任务，每项任务有明确的方案策略、量化指标与对应论文章节。表~\ref{tab:design-tasks} 给出任务—策略—指标的对照总览。

\begin{table}[htbp]
  \centering
  \small
  \caption{毕业设计任务分解}
  \label{tab:design-tasks}
  \begin{tabular}{>{\raggedright\arraybackslash}p{2.6cm}>{\raggedright\arraybackslash}p{4.0cm}>{\raggedright\arraybackslash}p{4.5cm}>{\raggedright\arraybackslash}p{2.5cm}}
    \toprule
    设计任务 & 方案/策略 & 考核指标 & 对应章节 \\
    \midrule
    T1: 弗兰卡 实时控制与 LeRobot 接口扩展 & 四层架构（实时层→网络层→抽象层→遥操作层）；实现弗兰卡机器人抽象与三维鼠标/同构主从两种遥操作入口；两阶段配对标定 & 控制频率 $\geqslant 30\,\mathrm{Hz}$；遥操作轨迹可复现性（标定后关节偏差 $<1^\circ$）；支持社区录制工具直接录制 & 第二章 \\
    \midrule
    T2: 桌面抓取示教数据采集与规范化 & LeRobot v2.1 格式；双相机（第三人称+腕部）；$10\,\mathrm{Hz}$ 帧率；任务语义标注 & $\geqslant 30$ 条有效轨迹；每条 $\geqslant 150$ 帧；字段模式定义 与 OpenPI 训练入口完全对齐 & 第二章 \\
    \midrule
    T3: $\pi_{0.5}$ 模型迁移与变体改进 & 全量微调基线 + 冻结骨干轻量化变体 + 短期记忆变体；观测/动作空间重映射；Flow Matching 训练 & 抓取成功率量化（10 次测试）；显存占用对比；失败恢复行为记录 & 第三章、第四章 \\
    \midrule
    T4: 实机闭环部署与在线执行策略对比 & 异步远程客户端 vs.\ 分段同步闭环；动作块截取执行（3--10 步/次）；高频重推理 & 同一权重下两种部署方式的成功率差异；闭环策略成功率目标 $\geqslant 90\%$ & 第三章 \\
    \bottomrule
  \end{tabular}
\end{table}

\noindent 设计任务详述如下。

T1：弗兰卡实时控制与 LeRobot 接口扩展。
设计目标是使 Franka FR3 机械臂被 LeRobot 框架视为标准机器人，从而可复用社区录制与训练工具链。方案上，首先在 RT PC 部署基于 libfranka 与 Ruckig 的实时服务进程，通过 TCP 套接字与上位机解耦；其次在 LeRobot 工厂中注册机器人抽象，将关节状态、夹爪开合度与多相机 RGB 映射为标准张量字典；再次实现三维鼠标笛卡尔增量遥操作（适用于快速粗定位）与同构主从关节空间遥操作（适用于精细接触阶段）两种示教入口；最后设计两阶段配对标定流程，自动解算七轴缩放偏置与夹爪量程并将结果版本化。考核指标：控制频率 $\geqslant 30\,\mathrm{Hz}$；标定后主从关节偏差均值 $<1^\circ$；支持社区录制工具直接录制。

T2：桌面抓取示教数据采集与规范化。
设计目标是以"抓取绿色苹果并放置至桌面左上角"为任务实例，采集符合 LeRobot v2.1 规范的桌面抓取数据集。方案上，采用第三人称固定相机与腕部手眼相机的双视角布局，以 $10\,\mathrm{Hz}$ 帧率同步录制，在接触敏感阶段通过同构主从臂完成精细示教；按 LeRobot v2.1 目录结构将任务语义与轨迹信息存入元数据目录，帧级表格以列式存储格式存放，视频轨道单独挂载；在元信息中明确定义 8 维状态与动作特征的类型与形状，确保与训练管线字段一致。考核指标：$\geqslant 30$ 条有效轨迹，总帧数 $\geqslant 5,000$；字段与 OpenPI 训练入口完全对齐，无需中间转换。

T3：$\pi_{0.5}$ 模型迁移与变体改进。
设计目标是将 $\pi_{0.5}$ 预训练权重迁移至桌面抓取任务，并针对小样本、接触敏感的特性设计模型变体。方案上，首先将双相机映射为全局与手眼两路视觉输入，8 维状态/动作与 LeRobot 标准字段对齐，动作块长度设为 32；其次加载预训练权重完成全量微调作为实验基线；进而利用参数路径过滤机制精确冻结 VLM 骨干而仅微调动作专家，目标是将训练显存由约 160 GB 降至约 60 GB 且成功率不降；同时基于空间-时间可分离注意力机制构造短期记忆变体，当前帧保留完整空间细节、历史帧压缩为少量摘要 token，目标成功率 $\geqslant 90\%$ 并具备失败恢复能力。两项改进详见第四章。考核指标：全量微调基线成功率 $\geqslant 75\%$；冻结骨干变体训练显存约 60 GB，成功率不低于基线；短期记忆变体成功率 $\geqslant 90\%$，并记录到自主失败恢复行为。

T4：实机闭环部署与在线执行策略对比。
设计目标是量化不同在线执行策略对同一 VLA 权重在接触敏感任务上的性能影响，验证部署策略与模型结构同等重要的假设。方案上，实现异步远端推理（执行完整动作块，偏长时开环）与分段同步闭环（每次仅执行一小段，约 $0.27\,\mathrm{s}$ 后重观测再推理）两种部署方式，使用同一模型权重在相同桌面场景下各进行 10 次抓取测试，记录成功率与失败模式。考核指标：分段同步闭环成功率 $\geqslant 90\%$；两种部署方式成功率差异 $\geqslant 50$ 个百分点，验证闭环策略的决定性作用。

\section{主要贡献与创新点}

综合上述设计任务与实验结果，本文的主要贡献与创新点归纳如下。

首先，弗兰卡--LeRobot 一体化接口实现与可版本化标定链路。将实时控制、LeRobot Robot 抽象与多种遥操作入口统一在同一数据录制语义下，降低与 $\pi_{0.5}$/OpenPI 训练栈的对接成本；通过两阶段配对标定将关节缩放偏置与夹爪量程写入可版本化 JSON 配置，显著提升示教轨迹在装配与零点误差条件下的可重复性（T1）。

其次，符合社区规范的小规模桌面抓取数据集。以自采桌面抓取数据集（30 条轨迹，5,512 帧，双相机 RGB + 8 维状态/动作）为实例，展示任务语义、双相机互补视角与帧级特征如何嵌入 LeRobot v2.1 格式，为小样本 VLA 微调提供可直接消费的数据基础（T2）。

再次，$\pi_{0.5}$ 向静平台桌面任务的迁移适配与模型结构改进。完成观测--动作空间重定义与训练配置适配；提出基于参数路径过滤的冻结骨干轻量化策略（训练显存由约 160 GB 降至约 60 GB，批次大小为 4，成功率 $90\%$）和基于空间-时间可分离注意力机制的短期记忆变体（成功率 $100\%$，具有自主失败恢复能力），两项改进通过配置驱动设计（T3）。

最后，在线闭环策略对接触敏感任务决定性作用的系统验证。在同一模型权重下，分段同步闭环将抓取成功率从约 $5\%$（异步开环）提升至接近 $100\%$（分段闭环），通过受控对比实验量化了部署策略对实机性能的影响量级，为 VLA 在类似接触富集任务中的部署提供了可操作的经验准则（T4）。
